% GENERATED FILE. Edit canonical lessons, then run npm run generate:books.
% canonical-source-hash: fnv1a64:c2bf29ff928669ea
% canonical-lessons: LA-C35-meridies

\chapter{Noon and Afternoon}
\label{ch:la-noon-and-afternoon}
\hlchaptermodality{35}

\begin{chapteropening}
\textbf{By the end of this chapter:} \emph{I can say midday in Latin, and use ante merīdiem and post merīdiem to place any hour of the day.}

Say merīdiēs, place an hour with ante merīdiem or post merīdiem, and say why Rome needed no separate afternoon noun.
\end{chapteropening}

\section[merīdiēs]{merīdiēs — the word you already know, one sound-change later}

\label{lesson:LA-C35-meridies}



\begin{quote}
You already have \emph{medius diēs}. Rome had no separate word for \textquotedblleft{}afternoon\textquotedblright{}; it split the whole day around this word.
\end{quote}

\begin{sounds}
\textbf{Merīdiēs} — \textquotedblleft{}\textbf{midday, noon}\textquotedblright{} — is the \textbf{same phrase} as \textbf{medius diēs} (Chapter 12), just \textbf{contracted}: an earlier \emph{medīdiēs} dropped its repeated \emph{d}-sound (dissimilation) to become \emph{merīdiēs}, a fifth-declension noun in its own right (accusative \emph{merīdiem}). It's the ultimate root of English \textbf{meridian}, via the Latin adjective \textbf{merīdiānus} (\textquotedblleft{}of midday\textquotedblright{}).
\end{sounds}

\begin{culture}
Here's the twist: Spanish has \emph{tarde}, French has \emph{après-midi}, German has \emph{Nachmittag} — a single word naming \textquotedblleft{}the afternoon\textquotedblright{} as its own block of time. \textbf{Latin doesn't.} Instead, Romans split the \textbf{entire day} around \emph{merīdiēs} into two large halves: \textbf{ante merīdiem} (\textquotedblleft{}\textbf{before noon}\textquotedblright{}) and \textbf{post merīdiem} (\textquotedblleft{}\textbf{after noon}\textquotedblright{}) — the direct ancestors of English's own \textbf{a.m.} and \textbf{p.m.}, used on every clock and phone today, often without anyone realizing they're speaking Latin.

This wasn't just a loose figure of speech — it carried real legal weight. Under the \textbf{Twelve Tables}, Rome's foundational law code, a civil case had to be argued \textbf{before noon}, with both parties present; if one party failed to appear, the magistrate could rule \textbf{against them, in absentia, once it passed noon}. \emph{Ante} and \emph{post merīdiem} divided not just clock time, but courtroom procedure.
\end{culture}

\begin{cousinweb}
You've already met \emph{tardus} (\textquotedblleft{}slow, late\textquotedblright{}) through Spanish's own \emph{tarde} — the very word Spanish reshaped into its \textquotedblleft{}afternoon.\textquotedblright{} In Latin itself, \emph{tardus} never made that leap: it stayed a plain adjective, \textquotedblleft{}slow, late,\textquotedblright{} and never became a Latin noun for \textquotedblleft{}afternoon.\textquotedblright{} Spanish's shift from \textquotedblleft{}late\textquotedblright{} to \textquotedblleft{}the afternoon\textquotedblright{} happened entirely on its own, after Latin.
\end{cousinweb}

\subsection*{Your turn}
Say these aloud:

\begin{itemize}
\raggedright
  \item \textquotedblleft{}merīdiēs\textquotedblright{} — the same medius diēs, contracted
  \item \textquotedblleft{}ante merīdiem,\textquotedblright{} \textquotedblleft{}post merīdiem\textquotedblright{} — before/after noon, the source of a.m./p.m.
  \item the honest twist — no separate Latin word for \textquotedblleft{}afternoon,\textquotedblright{} just a line drawn at noon
\end{itemize}

\begin{tcolorbox}[breakable,colback=teal!4,colframe=teal!35!black,title={Before you move on}]
What earlier phrase is \emph{merīdiēs} a contracted form of? (\emph{Medius diēs}, Chapter 12.) What two English abbreviations, used on every clock, come directly from \emph{ante merīdiem} and \emph{post merīdiem}? (\textbf{a.m.} and \textbf{p.m.}.) Did this noon-based division matter beyond just telling time? (\textbf{Yes} — under the Twelve Tables, Roman court cases had a real legal deadline at noon.) Does Latin have its own separate word for \textquotedblleft{}afternoon,\textquotedblright{} the way Spanish has \emph{tarde}? (\textbf{No} — Romans split the whole day around \emph{merīdiēs} instead.)
\end{tcolorbox}
